\documentclass[12pt, a4paper]{article}

% === Pacchetti ===
\usepackage[utf8]{inputenc}
\usepackage[T1]{fontenc}
\usepackage[italian]{babel}
\usepackage{geometry}
\usepackage{booktabs}
\usepackage{array}
\usepackage{xcolor}
\usepackage{colortbl}
\usepackage{enumitem}
\usepackage{fancyhdr}
\usepackage{titlesec}
\usepackage{microtype}
\usepackage{hyperref}
\usepackage{tcolorbox}
\usepackage{parskip}
\usepackage{lmodern}
\usepackage{amsmath}
\usepackage{amssymb}
\usepackage{tabularx}

% === Geometria pagina ===
\geometry{
  top=2.5cm, bottom=2.5cm,
  left=3cm,  right=2.5cm,
  headheight=15pt
}

% === Palette colori ===
\definecolor{primaryblue}{HTML}{1A3A5C}
\definecolor{accentblue}{HTML}{2E7EBA}
\definecolor{lightgray}{HTML}{F5F7FA}
\definecolor{medgray}{HTML}{AABCD0}
\definecolor{passgreen}{HTML}{2E7D32}
\definecolor{failred}{HTML}{C62828}
\definecolor{warnAmber}{HTML}{E65100}
\definecolor{roweven}{HTML}{EFF4F9}
\definecolor{sessionblue}{HTML}{EAF2FB}

% === Hyperref ===
\hypersetup{
  colorlinks=true,
  linkcolor=accentblue,
  urlcolor=accentblue,
  pdftitle={Report di Sessione --- 11 Giugno 2026},
  pdfauthor={Paolo},
  pdfsubject={Tesi Magistrale},
}

% === Header / Footer ===
\pagestyle{fancy}
\fancyhf{}
\fancyhead[L]{\small\color{medgray}\textit{Tesi Magistrale: GraphRAG per Molecular Tumor Board}}
\fancyhead[R]{\small\color{medgray}Giugno 2026}
\fancyfoot[C]{\small\color{medgray}\thepage}
\renewcommand{\headrulewidth}{0.4pt}
\renewcommand{\footrulewidth}{0pt}

% === Titoli sezione ===
\titleformat{\section}
{\Large\bfseries\color{primaryblue}}
{\thesection.}{0.6em}{}
[\titlerule]
\titleformat{\subsection}
{\normalsize\bfseries\color{accentblue}}
{\thesubsection.}{0.5em}{}
\titlespacing*{\section}{0pt}{1.4em}{0.6em}
\titlespacing*{\subsection}{0pt}{1em}{0.3em}

% === Box colorati ===
\tcbuselibrary{skins,breakable}
\newtcolorbox{summarybox}{
  enhanced, breakable,
  colback=lightgray, colframe=primaryblue,
  boxrule=1.2pt, arc=4pt,
  fonttitle=\bfseries\color{white},
  title={Panoramica della Sessione},
}
\newtcolorbox{insightbox}{
  enhanced, breakable,
  colback=sessionblue, colframe=accentblue,
  boxrule=1pt, arc=4pt,
  fonttitle=\bfseries\color{white},
  title={Implicazione Clinica},
}

% === Macro colori ===
\newcommand{\pass}[1]{{\color{passgreen}\textbf{#1}}}
\newcommand{\fail}[1]{{\color{failred}\textbf{#1}}}
\newcommand{\warn}[1]{{\color{warnAmber}\textbf{#1}}}
\newcommand{\bench}[1]{{\texttt{\textbf{#1}}}}

% ============================================================
\begin{document}

% === Titolo ===
\begin{center}
  {\LARGE\bfseries\color{primaryblue} Report di Sessione --- 11 Giugno 2026}\\[0.4em]
  {\large\color{accentblue} Sistema GraphRAG Agentico per Molecular Tumor Board}\\[0.2em]
  {\small\color{medgray}\textit{Tesi Magistrale --- Paolo}}
\end{center}
\vspace{0.5em}\hrule\vspace{1.2em}

% === Panoramica ===
\begin{summarybox}
Sessione di lavoro intensiva focalizzata su tre aree: aggiornamento dei risultati benchmark post-fix ESCAT, analisi e posizionamento rispetto a nuovi paper di letteratura, e preparazione all'incontro con il professore. La sessione ha prodotto risultati di benchmark definitivi, un report di analisi letteratura, una bozza di email al professore, e l'identificazione di un bug critico nella KB (nodo mancante per \bench{BENCH-004}).
\end{summarybox}

% === Sezione 1 ===
\section{Risultati Benchmark Aggiornati --- RQ1 Definitivi}

\subsection{Metriche Oggettive (post-fix ESCAT Interpreter)}

\begin{table}[htbp]
\centering
\small
\renewcommand{\arraystretch}{1.3}
\caption{Metriche oggettive di confronto per l'estrazione e grounding}
\begin{tabularx}{\textwidth}{l c c r}
\toprule
\rowcolor{primaryblue}
\color{white}\textbf{Metrica} & \color{white}\textbf{GraphRAG} & \color{white}\textbf{Zero-Shot} & \color{white}\textbf{Delta} \\
\midrule
Drug Anchor Check & 96.7\% (29/30) & 96.7\% (29/30) & 0.0\% \\
\rowcolor{roweven}
PMID Coverage Rate & 60.0\% & 0.0\% & \pass{+60.0\%} \\
PMID Hallucination Rate & 0.0\% & 100.0\% & \pass{-100.0\%} \\
\rowcolor{roweven}
ESCAT Tier Match & 100.0\% & 0.0\% & \pass{+100.0\%} \\
\bottomrule
\end{tabularx}
\end{table}

\subsection{Metriche Qualitative LLM-as-Judge}

\begin{table}[htbp]
\centering
\small
\renewcommand{\arraystretch}{1.3}
\caption{Valutazione qualitativa tramite LLM-as-Judge (Giudice: MiniMax-M2.5)}
\begin{tabularx}{\textwidth}{l c c r}
\toprule
\rowcolor{primaryblue}
\color{white}\textbf{Criterio} & \color{white}\textbf{GraphRAG} & \color{white}\textbf{Zero-Shot} & \color{white}\textbf{Delta} \\
\midrule
Score Totale & 3.72 & 3.92 & \fail{-0.20} \\
\rowcolor{roweven}
Completezza & 3.60 & 3.76 & \fail{-0.16} \\
Utilità Clinica & 3.53 & 3.67 & \fail{-0.13} \\
\rowcolor{roweven}
Fedeltà Evidenze & 3.90 & 4.13 & \fail{-0.23} \\
Accuratezza Clinica & 3.83 & 4.11 & \fail{-0.27} \\
\bottomrule
\end{tabularx}
\end{table}

\subsection{Lettura Corretta dei Risultati}
I risultati vanno letti su due livelli separati:
\begin{itemize}[leftmargin=1.5em, itemsep=3pt]
  \item \textbf{Livello 1 --- Metriche oggettive e verificabili:} GraphRAG domina nettamente. ESCAT al 100\%, zero hallucination sui PMID, coverage del 60\% vs 0\%. Queste metriche sono calcolabili automaticamente e non dipendono dal giudizio di un LLM.
  \item \textbf{Livello 2 --- LLM-as-judge:} Zero-shot appare superiore per ragioni di \textit{fluency bias} ampiamente documentato in letteratura. La prova del bias è interna: lo zero-shot allucina il 100\% dei PMID ma ottiene 4.13 su fedeltà evidenze --- un sistema che cita paper inesistenti non può avere fedeltà alle evidenze alta. Il giudice LLM valuta la fluidità percepita, non la correttezza fattuale. Questo è un finding metodologico originale della tesi, non un fallimento del sistema.
\end{itemize}

% === Sezione 2 ===
\section{Analisi Casi Drug Anchor Check --- Finding Narrativo Chiave}
Il Drug Anchor Check è pari a 96.7\% per entrambi i sistemi (29/30), ma i casi errati sono \textbf{diversi e asimmetrici} --- questo rappresenta il dato narrativamente più importante della sessione.

\subsection*{BENCH-004 --- GraphRAG sbaglia, Zero-Shot indovina}
\begin{itemize}[leftmargin=1.5em, itemsep=2pt]
  \item \textbf{Profilo:} \textit{ERBB2} Amplificazione, Carcinoma Mammario HER2+
  \item \textbf{Farmaco atteso:} Trastuzumab + Pertuzumab (combinazione)
  \item \textbf{Errore GraphRAG (score 0.5):} La Knowledge Base mancava del nodo per la combinazione completa. Si tratta di un problema di KB, non di pipeline. \textbf{Identificato come bug e risolto in sessione} tramite aggiunta del nodo e delle relazioni corrette.
  \item \textbf{Zero-Shot corretto:} La memoria parametrica dell'LLM riconosce lo standard storico consolidato del trattamento di prima linea per il tumore mammario HER2+.
  \item \textbf{Tipo di errore:} Tecnico, risolto con aggiornamento della KB.
\end{itemize}

\subsection*{BENCH-018 --- Zero-Shot sbaglia, GraphRAG indovina}
\begin{itemize}[leftmargin=1.5em, itemsep=2pt]
  \item \textbf{Profilo:} \textit{KRAS} G12C, Carcinoma del Colon-Retto (CRC)
  \item \textbf{Farmaco atteso:} Sotorasib + Panitumumab (combinazione)
  \item \textbf{Errore Zero-Shot (score 0.0):} Raccomanda la monoterapia con Sotorasib (standard per il tumore polmonare non a piccole cellule, NSCLC), ignorando che nel CRC l'inibizione del solo KRAS G12C causa una rapida riattivazione del feedback biologico mediato da EGFR, rendendo obbligatorio l'uso combinato di un anti-EGFR (Panitumumab). Approvazione FDA ad agosto 2023 sulla base del trial CodeBreaK 300 (NEJM 2023).
  \item \textbf{GraphRAG corretto:} Il Knowledge Graph contiene la relazione strutturata specifica per la combinazione KRAS G12C in CRC con Panitumumab, indipendente dalla distribuzione statistica del corpus di pre-training dell'LLM.
  \item \textbf{Tipo di errore:} Strutturale, non correggibile tramite prompt engineering. Si tratta del classico bias da distribuzione dei dati di training: la mutazione KRAS G12C associata a Sotorasib è un pattern statisticamente dominante per NSCLC, e lo Zero-Shot lo applica in modo errato anche al tumore del colon.
\end{itemize}

\begin{insightbox}
Il caso \bench{BENCH-018} illustra esattamente il rischio clinico dell'uso di LLM generici (\textit{vanilla}) in oncologia di precisione: un medico che si affidasse allo zero-shot per un paziente CRC con mutazione KRAS G12C riceverebbe una raccomandazione errata e potenzialmente dannosa. Il Knowledge Graph codifica la specificità istologica in modo deterministico e strutturato, non dipendente dalla frequenza statistica dei testi di addestramento. Questo rappresenta la giustificazione d'uso centrale della tesi.
\end{insightbox}

% === Sezione 3 ===
\section{Fix ESCAT Interpreter --- Completato}

\subsection{Problema identificato}
La mappatura ESCAT statica precedentemente adottata (\texttt{A $\rightarrow$ I-A, B $\rightarrow$ II-B}) produceva errori sistematici su tre categorie di alterazioni: varianti di resistenza, varianti tumor-agnostic e indicazioni off-label. La causa principale (\textit{root cause}) era la mancata verifica di corrispondenza tra l'istologia dell'evidenza nel grafo e il tipo di tumore del paziente.

\subsection{Fix implementato}
Sono stati introdotti quattro miglioramenti significativi nel file \texttt{variant\_interpreter.py}:

\textbf{1. Funzione \texttt{\_is\_same\_tumor()}} --- esegue un confronto fuzzy tra la patologia descritta nell'evidenza del grafo e il tumore effettivo del paziente:
\begin{tcolorbox}[colback=lightgray, colframe=primaryblue, title=Verifica corrispondenza istologica, boxrule=1pt, arc=4pt]
\begin{verbatim}
def _is_same_tumor(evidence_disease: str, patient_tumor: str) -> bool:
    if not evidence_disease or not patient_tumor:
        return False
    ev = evidence_disease.lower().strip()
    pt = patient_tumor.lower().strip()
    return ev == pt or ev in pt or pt in ev
\end{verbatim}
\end{tcolorbox}

\textbf{2. Arricchimento di \texttt{evidences\_str}} --- aggiunta del flag esplicito \texttt{Match istologico: SÌ/NO} passato come contesto all'agente decisore.

\textbf{3. Fallback intelligente} --- rimozione della mappatura rigida e introduzione di una logica contestuale basata sulla corrispondenza tumorale:
\begin{itemize}[leftmargin=1.5em, itemsep=2pt]
  \item Significatività \texttt{Resistance} $\rightarrow$ \texttt{non determinato}
  \item Stesso tumore: Livello A $\rightarrow$ \texttt{I-A}, Livello B $\rightarrow$ \texttt{I-B}, Livello C $\rightarrow$ \texttt{II-A}, Livello D $\rightarrow$ \texttt{II-B}
  \item Tumore diverso (off-label) con declassamento: Livello A/B $\rightarrow$ \texttt{II-B}, Livello C $\rightarrow$ \texttt{non determinato}
\end{itemize}

\textbf{4. Integrazione nel prompt} delle regole di interpretazione: se \texttt{Match istologico = NO}, declassa il tier; se si tratta di evidenze di resistenza, assegna sempre \texttt{non determinato}.

\subsection{Risultato}
La correttezza nell'assegnazione dei livelli di evidenza (\textit{ESCAT Tier Match}) è salita dal 73.3\% al \pass{100.0\%}.

% === Sezione 4 ===
\section{Nuovi Paper di Letteratura --- Analisi e Posizionamento}

\subsection{Paper acquisiti e analizzati}
\begin{itemize}[leftmargin=1.5em, itemsep=4pt]
  \item \textbf{Aldea et al. (2025) --- IASLC Consensus Statement on Molecular Tumor Boards}\\
  \textit{J Thorac Oncol. 2025;20:1594--1614}\\
  Documento clinico normativo che sancisce l'uso di scale di azionabilità (ESCAT/OncoKB) come standard per gli MTB. La frase chiave: \textit{``each ranking applies to a specific drug-target pair, not the alteration alone''} valida la scelta architetturale dell'ESCAT Interpreter basata su coppie strutturate. Da posizionare nell'Introduzione e nella Metodologia.
  
  \item \textbf{Berman et al. (2025) --- RAG\_MTB}\\
  \textit{JMIR Med Inform. 2025. PMID: 40053768}\\
  Sviluppo di una pipeline RAG testuale basata su PubMed per le raccomandazioni di terapia nel contesto MTB dell'Università di Tubinga. Rappresenta il confronto tecnologico diretto in letteratura: evidenzia i tre vantaggi strutturali del nostro lavoro (uso di KB curata rispetto a testi raw, controllo allucinazioni referenziali via relazione \texttt{CITED\_IN}, e classificazione ESCAT strutturata).
  
  \item \textbf{Loaiza-Bonilla et al. (2026) --- OncoMultiAgentKB}\\
  \textit{ESMO Real World Data and Digital Oncology. 2026. DOI: 10.1016/j.esmorw.2026.100706}\\
  Sistema multi-agente e knowledge graph oncologico proprietario valutato su 3.804 pazienti reali. Dimostra che la rimozione del grounding su grafo degrada l'F1-score da 0.82 a 0.79. Costituisce una validazione industriale esterna del nostro approccio GraphRAG. I nostri differenziatori principali includono la complessità del task (therapy recommendation vs trial matching), il vincolo della linea terapeutica e l'adozione di modelli open-source locali.
\end{itemize}

\subsection{Aggiornamento del Posizionamento}
Il posizionamento scientifico della tesi si articola ora in modo organico su quattro livelli:

\begin{table}[htbp]
\centering
\small
\renewcommand{\arraystretch}{1.3}
\caption{Quadro di posizionamento della tesi nello stato dell'arte}
\begin{tabularx}{\textwidth}{l l X}
\toprule
\rowcolor{primaryblue}
\color{white}\textbf{Livello} & \color{white}\textbf{Paper} & \color{white}\textbf{Funzione} \\
\midrule
Clinico & Aldea et al. (IASLC 2025) & ESCAT è lo standard clinico MTB riconosciuto. \\
\rowcolor{roweven}
Tecnico RAG & Berman et al. (JMIR 2025) & RAG testuale su PubMed ha limiti strutturali che la KB Neo4j risolve. \\
Tecnico multi-agente & Loaiza-Bonilla et al. (ESMO 2026) & KB oncologica + multi-agente è l'architettura vincente anche su scala industriale. \\
\rowcolor{roweven}
Contributo originale & Tesi (GraphRAG MTB) & Assegnazione ESCAT + vincolo linea + resistenza + open-source locale. \\
\bottomrule
\end{tabularx}
\end{table}

\subsection{Frase di Positioning per la Discussione con il Relatore}
\textit{``Il sistema proposto si colloca all'intersezione tra tre tendenze convergenti nella letteratura recente: la validazione clinica dell'uso di ESCAT come standard per le raccomandazioni MTB (Aldea et al., IASLC 2025), la dimostrazione che il RAG testuale su PubMed soffre di limitazioni strutturali nella gestione delle citazioni (Berman et al., JMIR 2025), e la validazione industriale su larga scala che la combinazione KB oncologica strutturata + agenti LLM specializzati riduce le allucinazioni e migliora la precision rispetto al solo LLM (Loaiza-Bonilla et al., ESMO 2026). La tesi contribuisce a questo panorama con una pipeline open-source completa che integra classificazione ESCAT, gestione della linea terapeutica e verifica strutturale delle citazioni --- elementi assenti o parziali in tutti e tre i sistemi di riferimento.''}

% === Sezione 6 ===
\section*{6. Ottimizzazioni Tecniche}
\begin{itemize}[leftmargin=1.5em, itemsep=3pt]
  \item \textbf{Driver Neo4j Persistente:} È stato rimosso il pattern inefficiente di creazione e chiusura del driver per ogni singola query, sostituendolo con un'unica istanza persistente e riutilizzabile. Il tempo totale per il calcolo delle metriche oggettive è sceso da diversi minuti a \pass{meno di 4 secondi}.
  \item \textbf{Aggiornamento KB per BENCH-004:} Aggiunto il nodo della terapia di combinazione e le relative relazioni con l'evidenza per \bench{BENCH-004} (\textit{ERBB2} Amplificazione, combinazione Trastuzumab + Pertuzumab, tumore mammario HER2+). Con questa integrazione, l'accuratezza terapeutica (Drug Anchor Check) di GraphRAG ha il potenziale di salire a 30/30 (100\%).
\end{itemize}

% === Sezione 7 ===
\section*{7. Stato Attuale del Progetto}

\subsection*{Completato in data odierna}
\begin{itemize}[leftmargin=1.8em, itemsep=2pt, label=\pass{[\textsf{x}]}]
  \item Fix logico dell'ESCAT Interpreter (precisione salita dal 73.3\% al 100\%)
  \item Benchmark RQ1 completato, stabile ed eseguito con successo su tutti i 30 casi
  \item Analisi di dettaglio e posizionamento scientifico di 3 nuovi paper clinico-tecnologici
  \item Identificazione e fix del bug della Knowledge Base relativo a \bench{BENCH-004}
  \item Stesura del report di analisi della letteratura (\texttt{report\_nuovi\_paper\_tesi.md})
  \item Bozza della mail di aggiornamento per il professore
  \item Analisi qualitativa ed evidenza del comportamento asimmetrico sui casi errati di Drug Anchor
\end{itemize}

\subsection*{Da pianificare per la prossima sessione}
\begin{itemize}[leftmargin=1.8em, itemsep=2pt, label={[\textsf{\phantom{x}}]}]
  \item Rieseguire il benchmark completo post-fix KB per confermare il Drug Anchor Check a 30/30 (100\%)
  \item Implementare la metrica di accuratezza citazionale automatica tramite interrogazione \texttt{CITED\_IN} nel grafo (confronto con MedGraphRAG)
  \item Calcolare il valore di precision@k approssimato per il Trial Matcher (confronto con TrialGPT)
  \item Aggiungere la discussione sulla linea terapeutica come fattore differenziatore nel report delle riflessioni
  \item Avviare la stesura della sezione ``Lavori Correlati'' nella tesi integrando gli 8 paper catalogati
  \item Concordare con il relatore la formulazione delle domande di ricerca rimanenti (RQ2, RQ3, RQ4)
\end{itemize}

\subsection*{Stato delle Research Questions (RQ)}
\begin{table}[htbp]
\centering
\small
\renewcommand{\arraystretch}{1.3}
\caption{Stato di avanzamento delle domande di ricerca}
\begin{tabularx}{\textwidth}{l X c}
\toprule
\rowcolor{primaryblue}
\color{white}\textbf{RQ} & \color{white}\textbf{Descrizione} & \color{white}\textbf{Stato} \\
\midrule
RQ1 & GraphRAG vs LLM Zero-Shot & \pass{Completato (Risultati Definitivi)} \\
\rowcolor{roweven}
RQ2 & Riproducibilità ESCAT + Accuratezza Citazionale & \warn{Parziale (Metriche da definire)} \\
RQ3 & Potenziale di automazione nella preparazione MTB & \texttt{Da iniziare} \\
\rowcolor{roweven}
RQ4 & Routing dinamico MoE vs Flusso sequenziale fisso & \texttt{Da iniziare} \\
\bottomrule
\end{tabularx}
\end{table}

% === Sezione 8 ===
\section*{8. Paper di Riferimento Completi (aggiornato)}

\begin{table}[htbp]
\centering
\scriptsize
\renewcommand{\arraystretch}{1.3}
\caption{Elenco sistematico dei paper di riferimento inclusi nella tesi}
\begin{tabularx}{\textwidth}{c p{4.2cm} c p{2.2cm} X}
\toprule
\rowcolor{primaryblue}
\color{white}\textbf{\#} & \color{white}\textbf{Paper} & \color{white}\textbf{Anno} & \color{white}\textbf{Venue} & \color{white}\textbf{Ruolo in Tesi} \\
\midrule
1 & Edge et al. --- GraphRAG & 2025 & Microsoft / arXiv & Fondamento teorico e logica GraphRAG \\
\rowcolor{roweven}
2 & Wu et al. --- MedGraphRAG & 2024 & Oxford / arXiv & Grounding strutturato su grafo in ambito clinico \\
3 & Kim et al. --- MDAgents & 2024 & MIT / NeurIPS & Orchestrazione multi-agente adattiva (RQ4) \\
\rowcolor{roweven}
4 & Jin et al. --- TrialGPT & 2024 & NIH / Nat. Comm. & Blueprint architetturale per Trial Matcher \\
5 & D\'iaz Cant\'on --- Agentic Shift & 2026 & Cancer Sci Res & Inquadramento clinico iniziale e motivazione \\
\rowcolor{roweven}
6 & Aldea et al. --- IASLC Consensus & 2025 & J. Thorac. Oncol. & Validazione dell'uso di scale ESCAT negli MTB \\
7 & Berman et al. --- RAG\_MTB & 2025 & JMIR Med Inform & Benchmark tecnologico diretto con RAG testuale \\
\rowcolor{roweven}
8 & Loaiza-Bonilla et al. --- OncoMultiAgentKB & 2026 & ESMO RWD & Validazione architettura su scala industriale \\
\bottomrule
\end{tabularx}
\end{table}

\vfill
\begin{center}
  {\color{medgray}\rule{0.5\linewidth}{0.4pt}}\\[0.3em]
  {\small\color{medgray}\textit{Report di sessione tradotto --- Sessione del 11 Giugno 2026}}
\end{center}

\end{document}
